\section{From local phases to minimal models}\label{sec:global}

We first isolate the local assertion and deduce the global theorem from it.
Throughout this section, until Proposition~\ref{prop:field-descent}, the base
field is $\C$.  Log discrepancies of divisorial valuations are homogeneous:
if $u=c\ord_E$ with $c>0$, then $A(u)=cA(\ord_E)$.  In particular,
``integral-valued'' does not require a divisorial valuation to be primitive.

At a centre where a volume form $\eta$ generates the canonical sheaf,
its form discrepancy $A_\eta$ is the ordinary boundary-zero discrepancy.
We use \eqref{glob:form-discrepancy} and its finite-extension compatibility
from Section~\ref{sec:conventions}, with the actual normalization of the
restricted valuation.

For a finite-order character, its angle belongs to $\R/\Z$, with the
convention that its value is $\exp(2\pi i\,\phaseangle(\chi))$.
Characters below describe the actions on functions and forms.

\begin{theorem}[Local phase assertion]\label{thm:phase}
Let $(V_i,o_i)$ be a sequence of pointed algebraic Gorenstein klt germs of a
fixed dimension $n\geq2$, with $o_i$ closed.  Suppose that a finite cyclic
group $G_i$ acts on an affine representative of each germ, fixes $o_i$, and
admits a semi-invariant canonical generator $\theta_i$, a volume form with
zero divisor near $o_i$.  Choose
an arbitrary element $\gamma_i\in G_i$ for each $i$.

After passage to a subsequence, there are divisorial valuations $w'_i$
centred at $o_i$, a constant $C>0$, and a positive integer $\ell$, independent
of $i$ on that subsequence, such that
\[
 A_{\theta_i}(w'_i)\leq C
\]
and, for every nonzero rational semi-invariant function $h$ on $V_i$ of
character $\chi$,
\begin{equation}\label{glob:phase-congruence}
 w'_i(h)\equiv\ell\,\phaseangle\bigl(\chi(\gamma_i)\bigr)
 \pmod{\Z}.
\end{equation}
In particular, $w'_i$ restricts to an integral-valued valuation of
$\C(V_i)^{G_i}$.
\end{theorem}

The proof occupies Sections~\ref{sec:cones}--\ref{sec:rounding}, using
Proposition~\ref{prop:integral-jump}, whose proof is supplied in
Section~\ref{sec:integral-jump} and the subsequent sections.  We now give all
steps of the global deduction.

\subsection{Termination of smooth canonical programs}

\begin{proposition}\label{prop:smooth-termination}
Assume Theorem~\ref{thm:phase}.  Let $X$ be a smooth projective complex
variety with pseudo-effective canonical divisor.  A $K_X$-MMP with scaling
of an effective ample rational divisor $H$, chosen so that $(X,H)$ is klt
and $K_X+H$ is ample, terminates with nef canonical divisor.
Consequently, on a smooth projective birational model $p:U\to X$ one has
\[
 p^*K_X=P+N,
 \qquad P\text{ nef and rational Cartier},\quad N\geq0
 \text{ rational Cartier}.
\]
\end{proposition}

\begin{proof}
The assertion is immediate in dimension zero; in dimension one a smooth
projective curve with pseudo-effective canonical divisor already has nef
canonical divisor.  Suppose that $n=\dim X\geq2$.  Such an $H$ is obtained
from sufficiently positive general divisors with small coefficients.  The
ordinary klt MMP, including existence of flips \cite{BCHM}, gives a program
\[
 X=X_0\dashrightarrow X_1\dashrightarrow X_2\dashrightarrow\cdots
\]
with $X_i$ projective, $\Q$-factorial and klt.  Write $H_i$ for the strict
transform of $H$ and $\lambda_i$ for the scaling number at its $i$th step.
If the program is infinite, the $\lambda_i$ lie in $(0,1]$, are
nonincreasing and rational: a contracted curve is $K_{X_i}$-negative and
$(K_{X_i}+\lambda_iH_i)$-trivial.  Only finitely many steps can be
divisorial, because each such step decreases the Picard number.

For a divisorial valuation $u$ of the common function field, set
\begin{equation}\label{glob:increments}
 A_i(u)=A_{X_i}(u),\qquad
 E_i(u)=A_i(u)-A_X(u),\qquad
 \Delta_i(u)=A_{i+1}(u)-A_i(u)\geq0.
\end{equation}
We use compatible canonical divisors on every common resolution.  At step
$j$ the pullbacks of the adjoints $K+\lambda_jH$ agree.  Their difference is
exceptional over the output and numerically trivial there, so this follows
from the negativity lemma.  The difference of the canonical pullbacks is
effective by the same lemma.  It follows, by summing the stepwise
identities, that for $0<t\leq\lambda_i$ the difference between the pullback
of $K_X+tH$ and that of $K_{X_i}+tH_i$ is effective and exceptional over
$X_i$, and its value at $u$ is
\begin{equation}\label{glob:fixed-difference}
 F_i(t;u)=\sum_{j<i}\left(1-\frac{t}{\lambda_j}\right)\Delta_j(u).
\end{equation}
Thus the transformed pair $(X_i,tH_i)$ is klt.  For rational $t$, this
effective exceptional difference identifies the systems of sufficiently
divisible multiples on the two sides after adding their fixed part.
In particular, $K_{X_i}+tH_i$ is big.  At $t=\lambda_i$ it is also nef,
and klt base point freeness makes it semiample.

We spell out why
\begin{equation}\label{glob:scaling-zero}
 \lambda_i\longrightarrow0.
\end{equation}
Otherwise choose $j$ after the last divisorial contraction, put $S=X_j$,
and choose a positive rational $\beta<\lim_i\lambda_i$.  The tail maps
from $S$ are small.  For $\beta\leq t\leq\lambda_j$, the boundaries $tH_j$
are klt and their adjoints are big.  Moreover, $H_j$ is big: sections of
sufficiently divisible multiples of $H$ push to sections of its transform.
Choose a decomposition
\[
 H_j\sim_\Q A_1+\cdots+A_r+G,
\]
where $G\geq0$ and the $A_a$ are general, small effective ample rational
divisors whose numerical classes span $N^1(S)_\R$.  In each $tH_j$ replace
only a fixed sufficiently small rational portion $\delta H_j$, with
$0<\delta<\beta$, by this expression.  The resulting boundaries remain
klt, since the entire compact interval can be checked on one log
resolution.  Small variations of the coefficients of the $A_a$ give a
rational polytope of klt boundaries with a common ample part and big
adjoints.  BCHM finiteness of marked ample models
\cite[Corollary~1.1.5]{BCHM} applies to this polytope.

Small $\Q$-factorial modifications identify the numerical divisor spaces.
For completeness, if a divisor is numerically trivial on one model, its
pullback and the pullback of its transform differ by a divisor exceptional
over the other model; that difference is numerically trivial over that
model and is zero by negativity.  Hence numerical equivalence is
transported in both directions.  At each scaling value $\lambda_i$, the
nef class $K_{X_i}+\lambda_iH_i$ can therefore be made ample by an
arbitrarily small rational perturbation in the displayed spanning
directions.  The perturbation may be chosen within the fixed polytope.
Smallness identifies the corresponding section systems, so the marked
model $S\dashrightarrow X_i$ is one of its finitely many ample models.
No marked model can recur: discrepancies are nondecreasing at each flip,
and some discrepancy increases strictly.  Indeed, otherwise the canonical
pullbacks would agree, contradicting the negative and positive
intersection signs on the two sides of the flip over its contraction
base.  This proves \eqref{glob:scaling-zero}.

For rational $0<t\leq1$, put $L_t=K_X+tH$.  Its asymptotic base order is
\[
 \sigma_u(t)=\inf_m\frac{1}{m}\,
 u\bigl(\mathfrak b(|mL_t|)\bigr),
 \qquad \sigma_u=\lim_{t\downarrow0}\sigma_u(t),
\]
where $m$ runs through sufficiently divisible positive integers.  The
values $\sigma_u(t)$ increase as $t$ decreases: a sufficiently divisible
multiple of $(t'-t)H$ is globally generated when $t'>t$.  At a scaling
value there is no asymptotic base order on $X_i$, by semiampleness, so
\eqref{glob:fixed-difference} gives
\begin{equation}\label{glob:sigma-squeeze}
 \sigma_u(\lambda_i)\leq E_i(u)\leq\sigma_u.
\end{equation}
The second inequality follows by fixing $i$, observing that
$\sigma_u(t)\geq F_i(t;u)$ for $0<t\leq\lambda_i$, and letting $t$ tend
to zero.

\subsubsection*{A uniform multiplier-ideal estimate}
For each positive integer $s$, consider
\[
 \mathcal J_s(t)=\mathcal J\bigl(X,s\|L_t\|\bigr).
\]
These ideals form a decreasing family as $t\downarrow0$, by the same
comparison of base ideals.  They stabilize to an ideal $\mathcal J_s$.
Here stabilization follows from global generation, not from a descending
chain condition for ideals.  Fix a very ample divisor $H_0$.  For each
$s$, choose a Cartier divisor $P_s$ so positive that
\[
 P_s-jH_0-K_X-sL_t\quad\text{is ample for every }0\leq t\leq1
 \text{ and }1\leq j\leq n.
\]
Asymptotic Nadel vanishing and Castelnuovo--Mumford regularity make all
$\mathcal J_s(t)\otimes\cO_X(P_s)$ globally generated
\cite[Corollary~1.5 and Theorem~1.8]{DEL}.  Their spaces of sections are nested subspaces of the
fixed finite-dimensional vector space $H^0(X,\cO_X(P_s))$, and therefore
eventually constant.  Generation then makes the ideals themselves
constant for all sufficiently small rational $t>0$.

For every divisorial $u$ one has
\begin{equation}\label{glob:multiplier-bounds}
 s\sigma_u(t)-A_X(u)
 \leq u(\mathcal J_s(t))\leq s\sigma_u(t).
\end{equation}
The first bound follows from the log-resolution formula for a divisible
level computing the asymptotic multiplier ideal.  For the second,
asymptotic subadditivity on the smooth variety $X$
\cite[Variants~2.4--2.5]{DEL} gives, for sufficiently divisible $m$,
\[
 \mathfrak b(|msL_t|)
 \subseteq\mathcal J\bigl(X,ms\|L_t\|\bigr)
 \subseteq\mathcal J_s(t)^m.
\]
Taking orders, dividing by $m$, and taking the asymptotic infimum proves
the bound.  Combining stabilization, \eqref{glob:scaling-zero},
\eqref{glob:sigma-squeeze}, and \eqref{glob:multiplier-bounds}, we obtain:
for every fixed positive integer $s$, there is an index $i(s)$ such that
for every $i\geq i(s)$ and every divisorial valuation $u$,
\begin{equation}\label{glob:uniform-squeeze}
 \frac{u(\mathcal J_s)}{s}
 \leq E_i(u)\leq\sigma_u
 \leq\frac{u(\mathcal J_s)+A_X(u)}{s}.
\end{equation}
The index $i(s)$ is independent of $u$.  In particular,
\begin{equation}\label{glob:small-increment}
 0\leq\Delta_i(u)\leq\frac{A_X(u)}s
 \qquad(i\geq i(s)).
\end{equation}

\subsubsection*{Countability of finite limiting changes}
We claim that there is a fixed countable subset $\mathscr E\subseteq\R$
containing every finite limit of $E_{i_a}(u_a)$, where $i_a\to\infty$,
the $u_a$ are integral-valued divisorial valuations, and
$A_X(u_a)\leq C$ for some constant $C$ depending on the sequence.

Choose a countable algebraically closed field $F\subset\C$ over which $X$
and all the countably many ideals $\mathcal J_s$ are defined.  For every
$f\in F(X_F)^*$, its zero and pole divisors on $X$ have positive log
canonical thresholds when nonzero.  Their orders at $u_a$ are therefore
bounded by constants times $A_X(u_a)$, so $u_a(f)$ is bounded.  Since these
values are integral, diagonal extraction over the countable field makes
them eventually constant for every $f$.  The eventual values define an
integral valuation $u_\infty$ on $F(X_F)$, trivial on $F$.

We next show that $u_\infty$ is either trivial or divisorial.  If it is
nontrivial and its centre is not a divisor, blow up the closure of its
centre and continue.  At each stage we work near the generic point of the
centre, where the ambient variety and that centre can be taken smooth.
The next centre lies above this generic smooth neighbourhood, where
the blowup is smooth. Taking a projective model and resolving outside
it therefore does not change the valuation calculation.  A centre of codimension at least two has an ideal of
positive integral value.  Blowing it up thus increases the value of the
relative Jacobian over $X_F$ by at least one.  If $N$ such blowups were
possible, choose a smooth affine neighbourhood $U=\Spec A$ of the last
centre on which this relative Jacobian divisor is principal, with equation
$b$.  Then $u_\infty(b)\geq N$.

Choose finitely many $F$-algebra generators of $A$, including the inverses
used to obtain this affine neighbourhood.  Their $u_a$-values are
eventually equal to their nonnegative $u_\infty$-values.  Because each
$u_a$ is trivial on $\C$, every polynomial in these generators with
complex coefficients has nonnegative value.  Thus $u_a$ has a centre on
$U_\C$ for all sufficiently large $a$.  Its centre need not be the base
change of the centre of $u_\infty$.  The Jacobian identity holds throughout
this chart and nevertheless gives
\[
 A_X(u_a)=A_{U_\C}(u_a)+u_a(b)\geq u_a(b)
 =u_\infty(b)\geq N.
\]
Taking $N>C$ is impossible.  The centre therefore eventually has
codimension one.  The valuation then dominates its discrete valuation
ring, is zero on its units, and is a positive integral multiple of that
DVR order.

There are only countably many projective models of $F(X_F)$ defined over
the countable field $F$, hence only countably many prime divisors on such
models and their integral multiples.  Adding the trivial valuation still
gives a countable set of possibilities for $u_\infty$.  For each fixed $s$,
the preceding finite-chart argument applied to finitely many local
generators of $(\mathcal J_s)_F$ gives
\[
 u_a(\mathcal J_s)=u_\infty((\mathcal J_s)_F)=:a_s
 \quad\text{for all sufficiently large }a.
\]
If $E_{i_a}(u_a)\to e$, \eqref{glob:uniform-squeeze} therefore implies
\begin{equation}\label{glob:unique-limit}
 \frac{a_s}{s}\leq e\leq\frac{a_s+C}{s}
 \qquad\text{for every positive integer }s.
\end{equation}
One limiting valuation determines at most one such $e$: a second sequence
may have a different bound $C'$, but the two resulting limits differ by at
most $\max\{C,C'\}/s$ for every $s$.  If $u_\infty$ is trivial, all $a_s$
are zero and $e=0$.  This proves the claimed countability, including its
independence from any subsequent choice of phases.

\subsubsection*{A common canonical index}
We claim that there is a positive integer $R$ such that
\begin{equation}\label{glob:common-index}
 RK_{X_i}\text{ is Cartier for every }i.
\end{equation}
Otherwise, after passage to a subsequence, choose closed points
$x_i\in X_i$ whose local canonical indices $r_i$ tend to infinity.  If the
indices were bounded, their least common multiple would give
\eqref{glob:common-index}.

We recall the index-one cover at such a point.  Choose a rational top form
$\eta$, write $K=\operatorname{div}(\eta)$, shrink about the point so that
$rK=\operatorname{div}(g)$, and normalize after adjoining $z$ with
$z^r=g$.  Minimality of $r$ and Kummer theory make this an extension of
degree $r$.  It is unramified in codimension one; ramification of
discrepancies shows that it is klt.  The form
\[
 \theta=\frac{\pi^*\eta}{z}
\]
has zero divisor and has primitive character under the cyclic cover group.
The cover is Gorenstein: its canonical sheaf is generated by $\theta$, and
klt singularities are Cohen--Macaulay.

There is exactly one point over the specified base point.  Indeed, a
finite group quotient has a transitive action on the points of each
geometric fibre.  If this orbit had more than one point, a nonconstant
character of its cyclic permutation representation would give a regular
semi-invariant $a$, of character $z^j$ with $0<j<r$, nonzero at every point
of the fibre.  One obtains $a$ by lifting its values on the reduced fibre
and projecting to the character space.  Its invariant power $a^r$ is
nonzero at the base point, so $a$ is a unit near the whole fibre.  The
function $a/z^j$ is invariant, and its divisor downstairs is $-jK$.
Thus $jK$ is principal near the point, contradicting the minimality of
$r$.  The unique point is fixed by the cover group, as required in
Theorem~\ref{thm:phase}.

Now fix any $\alpha\in\R/\Z$.  Because $r_i\to\infty$ and the character of
$\theta_i$ is primitive, choose group elements $\gamma_i$ so that the
angles of the $\theta_i$-characters at $\gamma_i$ tend to $\alpha$.
Apply Theorem~\ref{thm:phase}, and let $u_i$ be the restrictions of the
resulting $w'_i$ to the common downstairs field.  These restrictions are
integral-valued and divisorial, and quasi-\'etale discrepancy compatibility
gives
\[
 0<A_X(u_i)\leq A_i(u_i)=A_{\theta_i}(w'_i)\leq C.
\]
Choose a rational top form $\eta_i^0$ downstairs that generates the
canonical sheaf at the centre of $u_i$ on the smooth variety $X$.  With
pullback of forms understood, \eqref{glob:form-discrepancy} gives
\begin{equation}\label{glob:phase-change}
 E_i(u_i)=w'_i(\theta_i/\eta_i^0).
\end{equation}
The ratio is a rational semi-invariant with the character of $\theta_i$.
Consequently, after passage to the subsequence in
Theorem~\ref{thm:phase}, these nonnegative bounded changes have residues
tending to $\ell\alpha$.  Pass further so that they converge in $\R$.
Their limit belongs to the fixed countable set $\mathscr E$, and hence
$\ell\alpha$ belongs to its image $\overline{\mathscr E}$ in $\R/\Z$.
But
\[
 \bigcup_{\ell\geq1}\{\alpha\in\R/\Z:
                  \ell\alpha\in\overline{\mathscr E}\}
\]
is countable, whereas $\alpha$ was arbitrary.  The possible dependence of
$\ell$ on the subsequence does not change this contradiction.  This proves
\eqref{glob:common-index}.

\subsubsection*{Vanishing on late flips}
For every fixed positive integer $m$, all sufficiently late flips, with
negative contraction $f_i:X_i\to Z_i$, satisfy
\begin{equation}\label{glob:late-vanishing}
 R^k f_{i*}\cO_{X_i}(mK_{X_i})=0\qquad(k>0).
\end{equation}
To prove this, take a common smooth resolution
$p:U\to X_i$, $q:U\to X_{i+1}$ over $Z_i$, sufficiently high that the
fractional parts to be rounded have simple normal crossing support, and
put
\[
 \mathcal L_m=\cO_U\left(K_U+\left\lceil
                    (m-1)q^*K_{X_{i+1}}\right\rceil\right).
\]
The rational divisor $(m-1)q^*K_{X_{i+1}}$ is nef over both $X_i$ and
$Z_i$: a curve contracted by $p$ is also contracted over $Z_i$, and
$K_{X_{i+1}}$ is ample over $Z_i$.  It is big relative to both bases since
the relevant morphisms are birational; this includes $m=1$, whose divisor
is zero and whose generic fibres are zero-dimensional.  Relative
Kawamata--Viehweg vanishing \cite{KM98} therefore gives
\[
 R^kp_*\mathcal L_m=R^k(f_ip)_*\mathcal L_m=0\qquad(k>0).
\]

We claim that $p_*\mathcal L_m=\cO_{X_i}(mK_{X_i})$, where the right side
denotes the divisorial sheaf if $mK_{X_i}$ is not Cartier.  One inclusion
is checked in codimension one, where the flip and the resolution are
isomorphisms.  For the reverse inclusion, let a local rational section
$h$ of $\cO_{X_i}(mK_{X_i})$ be given.  The divisor
\[
 \operatorname{div}_U(h)+mp^*K_{X_i}
\]
is effective, because $\operatorname{div}_{X_i}(h)+mK_{X_i}$ is effective
and rational Cartier.  Before rounding, the difference between the
divisor defining $\mathcal L_m$ and $mp^*K_{X_i}$ has coefficient, at any
prime divisor of $U$,
\[
 A_i-1-(m-1)\Delta_i.
\]
For $i\geq i(m)$, \eqref{glob:small-increment} implies
\begin{equation}\label{glob:rounding-strict}
 A_i-1-(m-1)\Delta_i
 \geq -1+\frac{A_X}{m}>-1.
\end{equation}
Adding this to the effective pullback of the section divisor and taking
round-up leaves every integral coefficient nonnegative.  Hence $h$
belongs to $p_*\mathcal L_m$.  The pushforward identity and Leray prove
\eqref{glob:late-vanishing}.

Fix $R$ as in \eqref{glob:common-index}.  For each late step choose a
sufficiently positive very ample divisor $J$ on $Z_i$ so that the Cartier
divisor
\[
 D=-RK_{X_i}+f_i^*J
\]
and $D-K_{X_i}$ are ample.  This is possible by relative anti-ampleness
of $K_{X_i}$.  Effective base point freeness, applied to the klt pair
$(X_i,0)$ and the nef Cartier divisor $D$ with $D-K_{X_i}$ nef and big,
gives an integer $N=N(n)$ such that $L=ND$ is globally generated
\cite[Theorem~1.1 and Remark~1.2]{FujinoEffective}.  It is also ample.
Only now choose a stage late enough for \eqref{glob:late-vanishing} to hold
for the finite list
\[
 m=R+jNR,\qquad1\leq j\leq n.
\]
The integer $N$ is independent of the step and of $J$, so this order of
choices is legitimate.  For each such step, projection formula, Leray and
Serre vanishing on $Z_i$ give, for all sufficiently large integers $a$,
\begin{equation}\label{glob:regularity-vanishing}
 H^j\left(X_i,\cO_{X_i}
       (RK_{X_i}+af_i^*J-jL)\right)=0\qquad(j>0).
\end{equation}
Indeed the displayed divisor equals
$(R+jNR)K_{X_i}+(a-jN)f_i^*J$.

Castelnuovo--Mumford regularity with respect to the ample globally generated
$L$ now makes $\cO_{X_i}(RK_{X_i}+af_i^*J)$ globally generated.  To use the
usual very ample formulation, the morphism defined by $L$ is finite onto
its image, since $L$ is ample; push the sheaf to projective space, apply
regularity there using \eqref{glob:regularity-vanishing}, and pull back
the generating sections.  A globally generated line bundle has
nonnegative degree on every curve, contradicting its degree
$RK_{X_i}\cdot C<0$ on a curve contracted by $f_i$.  The hypothetical
infinite program is impossible.

Pseudo-effectivity of the canonical class persists at every finite step:
push forward the big systems of $K_X+tH$ for arbitrarily small rational
$t>0$, and take their limiting numerical classes.  It rules out a final
negative fibre contraction, since the restriction to a general
positive-dimensional fibre would be both pseudo-effective and
anti-ample.  The terminal model $Y$ therefore has nef canonical divisor.
On a common smooth projective resolution $p:U\to X$, $q:U\to Y$, the MMP
inequalities give
\[
 p^*K_X=q^*K_Y+N,\qquad N\geq0.
\]
Taking $P=q^*K_Y$ proves the last assertion.
\end{proof}

\subsection{Relative weak Zariski decompositions}

An NQC weak Zariski decomposition of a rational Cartier divisor $D$ over
$Z$ means a projective birational model $p:Y\to X$ and a numerical
decomposition $p^*D\equiv_Z P+N$, where $N\geq0$ and $P$ is a nonnegative
real linear combination of rational Cartier divisors nef over $Z$.
The decompositions constructed here have $P$ and $N$ rational Cartier,
so in particular their nef parts are NQC.

\begin{proposition}\label{prop:relative-wzd}
Assume Theorem~\ref{thm:phase}.  Let $X$ be a smooth complex
quasi-projective variety, projective over a normal quasi-projective
variety $Z$.  If $K_X$ is pseudo-effective over $Z$, it admits an NQC weak
Zariski decomposition over $Z$ with rational Cartier parts.
\end{proposition}

\begin{proof}
Replace $Z$ by the Stein base of the morphism, which is normal and finite
over its original image.  The morphism $f:X\to Z$ is then surjective with
connected fibres.  This replacement changes neither the contracted
curves nor relative numerical equivalence, nefness or pseudo-effectivity.
A very general fibre is smooth and projective by generic smoothness and
has pseudo-effective canonical divisor.  To see the latter assertion,
take a countable sequence of relative effective approximations to
$[K_X]$ and restrict to fibres outside their countably many bad loci.
Relative numerical equivalence restricts to numerical equivalence on
such a fibre, and $K_X$ restricts to its canonical divisor.  By BDPP,
these very general smooth fibres are non-uniruled
\cite[Theorem~0.2 and Corollary~0.3]{BDPP}.

Compactify the projective morphism and resolve away from $X$ to obtain a
projective morphism
\[
 \overline f:\overline X\longrightarrow\overline Z
\]
with $\overline X$ smooth projective, $\overline Z$ normal projective, and
the original morphism unchanged over $Z$.  If $\dim Z=0$, the variety $X$
is already projective and Proposition~\ref{prop:smooth-termination}
applies.  We may therefore assume $\dim Z>0$.

Let $r:\widetilde Z\to\overline Z$ be a smooth projective resolution.
Choose a sufficiently high multiple $H$ of a very ample divisor on
$\overline Z$, and a general branch divisor $B\in|2H|$.  The pulled-back
linear systems on both $\widetilde Z$ and $\overline X$ are base point
free.  Bertini therefore makes $r^*B$ and $\overline f^*B$ smooth,
nonempty reduced divisors.  The corresponding double covers
\[
 Z^\sharp\longrightarrow\widetilde Z,
 \qquad \pi:X^\sharp\longrightarrow\overline X
\]
are smooth integral projective varieties.  Integrality follows because a
nonempty reduced branch divisor cannot be the divisor of a square.  The
double-cover formula gives
\begin{equation}\label{glob:double-cover}
 \begin{aligned}
 K_{Z^\sharp}&\sim_\Q
     \pi_Z^*(K_{\widetilde Z}+r^*H),\\
 K_{X^\sharp}&\sim_\Q
     \pi^*(K_{\overline X}+\overline f^*H).
 \end{aligned}
\end{equation}
For sufficiently positive $H$, the first canonical divisor is big.
The rational map $X^\sharp\dashrightarrow Z^\sharp$ has the same very
general fibres as $f$.  The base is non-uniruled because its canonical
divisor is big, and those fibres are non-uniruled by the preceding
paragraph.  Hence $X^\sharp$ is non-uniruled: a covering family of rational
curves would either give a covering family on the base or consist of
vertical curves covering very general fibres.  BDPP now makes
$K_{X^\sharp}$ pseudo-effective.

Apply Proposition~\ref{prop:smooth-termination} to $X^\sharp$.  On a
projective birational model its canonical pullback is a sum of a nef
rational divisor and an effective rational divisor.  By
\eqref{glob:double-cover}, the same is true, after changing the nef
divisor by a rational principal divisor if necessary, for the pullback
of $\pi^*(K_{\overline X}+\overline f^*H)$.  Take a common equivariant
resolution of this model and its conjugate under the covering involution.
Writing $g:T\to X^\sharp$ for the resulting morphism and $\iota$ for the
involution, average the decomposition and its conjugate.  We obtain
\begin{equation}\label{glob:averaged-decomposition}
 g^*\pi^*(K_{\overline X}+\overline f^*H)
   =\overline P+\overline N,
 \quad \iota^*\overline P=\overline P,
 \quad \iota^*\overline N=\overline N,
\end{equation}
where $\overline P$ is nef and $\overline N\geq0$, both rational Cartier.

Let $q:T\to Y=T/\langle\iota\rangle$ be the finite quotient.  The normal
projective variety $Y$ is birational to $\overline X$.  An invariant
rational divisor descends as a rational Weil divisor by dividing its
coefficient at an upstairs prime by the corresponding ramification
index.  If the invariant divisor is rational Cartier, its descended
divisor is rational Cartier as well.  Indeed, after clearing
denominators, a Cartier divisor has one local equation on a semilocal
neighbourhood of the finite orbit above a downstairs point.  The product
of its translates is an invariant equation whose divisor is the group
order times the original divisor.  It descends to a rational local
equation for a multiple of the downstairs divisor.  This is the finite
norm argument, and applies at every point of $Y$.

Consequently $\overline P=q^*P$ and $\overline N=q^*N$ for rational
Cartier divisors $P,N$ on $Y$.  Nefness of $P$ follows by lifting any
curve to $T$ and dividing its intersection by the positive covering
degree; effectivity of $N$ is immediate from its coefficients.  Finite
pullback is injective on divisors, so
\eqref{glob:averaged-decomposition} descends to a decomposition of the
pullback of $K_{\overline X}+\overline f^*H$.  Restrict to the inverse
image of $X$.  The additional term pulled from $Z$ is numerically trivial
over $Z$, and we obtain the required relative weak Zariski decomposition
of $K_X$.
\end{proof}

\subsection{The generalized pair and its exact output}

We apply the results of Tsakanikas--Xie in their relative setting
\cite{TsakanikasXie}.  Their Theorem~D, stated as Theorem~5.2, says that
relative minimal-model existence for smooth varieties in dimension
$n-1$ upgrades an NQC weak Zariski decomposition of an $n$-dimensional
NQC generalized lc pair to a minimal model.  Starting with dimension
zero, Proposition~\ref{prop:relative-wzd} and induction therefore give
relative minimal models for all pseudo-effective smooth varieties in
every dimension.  Their Theorem~B, stated as Theorem~5.4, now gives
minimal models for all pseudo-effective NQC generalized lc pairs.
Their Theorem~2.7 identifies existence in this sense with existence in
the Birkar--Shokurov sense.

For the generalized rational pair in Theorem~\ref{thm:main}, the nef data
are NQC: a rational nef Cartier divisor, after a positive multiple, is
itself one of the permitted nef Cartier summands.  If its adjoint is
pseudo-effective, the preceding paragraph supplies a Birkar--Shokurov
minimal model.  If it is not pseudo-effective, the alternative hypothesis
of Tsakanikas--Xie's Theorem~A applies directly.  In either case, that
theorem, stated as Theorem~4.2, gives a terminating MMP with scaling
starting on the original $X$, without a $\Q$-factoriality assumption.
To meet its scaling hypothesis, choose an effective sufficiently positive
ample rational divisor $A$ such that
\[
 (X,B+A+M)\text{ is generalized lc},
 \qquad K_X+B+A+M\text{ is nef}.
\]
Such an $A$ is a general very ample member divided by a sufficiently
large integer.  On a fixed log resolution carrying the nef data, its
pullback is general in a base point free system and transverse to the
log boundary.  Small coefficients preserve the condition that all
boundary coefficients on this resolution are at most one.  Its ample
class can nevertheless be chosen large enough to make the adjoint nef.
This checks all hypotheses of the terminating-MMP theorem.

Let $X\dashrightarrow Y$ be the resulting finite birational program.
The boundary at each stage is the strict pushforward of $B$, and the nef
b-divisor is unchanged.  Every adjoint is rational Cartier.  Indeed, the
construction gives a rational Weil divisor which is real Cartier, and a
rational Weil divisor in the real span of Cartier divisors lies in their
rational span, by solving the finite system of rational linear equations
for its coefficients.

We verify the precise discrepancy inequalities rather than only a
numerical minimal-model condition.  A birational step is given by a
negative extremal contraction $S\to T$ and a positive model
$S^+\to T$, whose morphism is small; in the ordinary divisorial case the
positive model is $T$ itself.  This description also holds without
$\Q$-factoriality \cite[Remark~2.15]{TsakanikasXie}.  Denote the adjoints
by $D$ and $D^+$.  On a common resolution $p:U\to S$, $q:U\to S^+$ over
$T$, taken to carry the fixed nef data, put
\begin{equation}\label{glob:step-difference}
 F=p^*D-q^*D^+.
\end{equation}
It is $q$-exceptional.  For a $q$-contracted curve, its intersection is
that of $p^*D$, which is nonpositive because $D$ is negative on the
contracted extremal ray.  Thus $F$ is anti-nef over $q$, and negativity
gives $F\geq0$.

The coefficient of $F$ is strictly positive along the transform of each
prime divisor on $S$ contracted in the step.  If it were zero along such
a prime, choose a general point of its transform outside $\Supp F$
where $p$ is an isomorphism. Take a curve $C$ through this point in its
positive-dimensional $q$-fibre. It maps to a curve on $S$ and
is not contained in $\Supp F$. Effectivity therefore gives
$F\cdot C\geq0$, whereas $q^*D^+\cdot C=0$ and
$p^*D\cdot C<0$, a contradiction.  Since the same nef divisor occurs on
$U$ for both generalized pairs, \eqref{glob:step-difference} gives,
for every divisorial valuation $u$,
\begin{equation}\label{glob:generalized-monotonicity}
 A_{S,B_S+M_S}(u)\leq A_{S^+,B_{S^+}+M_{S^+}}(u).
\end{equation}
For an original prime contracted during the program, strictness holds at
the first step that contracts it, and persists thereafter.  Composition
therefore gives exactly the unscaled inequalities in
Theorem~\ref{thm:main}.  Smallness of the positive morphisms makes the
inverse of $X\dashrightarrow Y$ non-extractive.  Generalized lc is
preserved by \eqref{glob:generalized-monotonicity}.

The endpoint has nef adjoint, or has a final fibre-type extremal
contraction $h:Y\to T$, with connected fibres and normal projective
target of smaller dimension, on which $-D_Y$ is relatively ample.
We also justify its numerical rank without a factoriality assumption.
All vertical curves span its contracted ray.  If a Cartier divisor $L$
has degree zero on that ray, $L-D_Y$ is relatively ample.  Choose a
sufficiently positive ample divisor $H$ on $T$ and an effective ample
rational divisor
\[
 A\sim_\Q L+h^*H-D_Y
\]
with small general coefficients, so that adding $A$ to the boundary
preserves generalized lc as above.  The new adjoint
$D_Y+A\sim_\Q L+h^*H$ is nef, and hence pseudo-effective, over $T$.
Tsakanikas--Xie's Theorem~F, stated as Theorem~5.26, applies: its added
divisor is effective, real Cartier and relatively ample.  It gives an
MMP to a good minimal model.  Since the adjoint is already nef, that MMP
has no negative step, so $D_Y+A$ is semiample over $T$.

A relatively semiample divisor which is numerically trivial over $T$
is numerically a pullback from $T$.  In fact, its associated morphism
over $T$ sends each connected projective fibre of $h$ to a point: a
positive-dimensional image would contain a curve with positive degree
for the ample divisor on the image.  Its normal image is consequently
finite and birational over $T$, and equals $T$ by normality and the
contraction property.  Applied here, this shows that $L$, up to its
already specified pullback twist, is numerically from $T$.  A contracted
curve defines a rational hyperplane in $N^1(Y)_\R$; that hyperplane is
therefore precisely $h^*N^1(T)_\R$.  If semiampleness is expressed as a
positive real sum of semiample rational Cartier divisors, every summand
has nonnegative degree on a vertical curve; the vanishing of their sum
forces each degree to vanish, and the same pullback argument applies to
each summand.  The hyperplane has codimension one, giving
\[
 \rho(Y/T)=1,\qquad \rho(Y)-\rho(T)=1.
\]
The sign of the original adjoint determines the endpoint. Birational
steps preserve pseudo-effectivity of the adjoint: on a common
resolution the difference is effective and exceptional as in
\eqref{glob:step-difference}, and pushforward gives the reverse
implication. A pseudo-effective adjoint cannot be anti-ample on the
general positive-dimensional fibre of a contraction. Conversely, a nef
endpoint and the effective pullback differences make the original
adjoint pseudo-effective. Hence the pseudo-effective case ends with a
nef model and the other case with a Mori fibre space. Thus all
conclusions of Theorem~\ref{thm:main} hold over $\C$.

\subsection{Descent to arbitrary characteristic-zero fields}

\begin{proposition}\label{prop:field-descent}
The conclusion of Theorem~\ref{thm:main} over $\C$ implies its conclusion
over every algebraically closed field of characteristic zero.
\end{proposition}

\begin{proof}
Let $k$ be such a field.  Descend the given projective varieties, rational
divisors and nef data, together with a log resolution carrying those
data, to a countable algebraically closed subfield $F\subset k$.
Embed $F$ into $\C$.  The initial properties hold over $F$ and after
extension to $\C$: generalized lc is checked on the chosen log
resolution, and nefness is preserved under algebraically closed field
extension by spreading and specializing curves.  Apply the complex
result.  We will descend its finite MMP to $F$ step by step; the chosen
complex scaling divisor need not descend.

Suppose the current variety $S$ and adjoint $D$ are defined over $F$, and
their complex extensions have been identified with the corresponding
stage of the program.  Consider its next complex contraction
$f_\C:S_\C\to T_\C$.  Pull back a very ample line bundle from $T_\C$
and choose generating sections.  This bundle and its sections spread on
$S\times_F B$, for a nonempty integral finite-type $F$-scheme $B$.
After shrinking $B$, they generate the bundle on every fibre: failure
of generation has closed image under the projection, since $S$ is
projective.  Specialize at an $F$-point and take the Stein contraction
of the resulting map.  This gives $f:S\to T$ over $F$, with normal target
and connected fibres.

After extension to $\C$, the specialized line bundle is numerically
equivalent to the original one.  Indeed, $B$ is geometrically integral
because $F$ is algebraically closed.  On its connected complex base
change, the degree of the universal bundle restricted to any fixed
projective curve on $S_\C$ is constant.  Hence the two contractions
contract exactly the same curves.

We explain why this identifies their targets, including in the
fibre-type case.  Each morphism is constant on every connected projective
fibre of the other: a positive-dimensional image of a fibre component
contains a curve, and hyperplane sections of its projective preimage
give a curve mapping onto it, contrary to equality of the contracted
curves.  The constants agree on intersecting components and hence on
the whole connected fibre.  The reduced joint image in the product of
the two targets therefore has finite projections to both.  After
normalization, these projections are birational: the function field of
a normal contraction target is algebraically closed in $\C(S)$, so
there is no nontrivial finite intermediate field.  Both targets are
normal, and finite birational morphisms to them are isomorphisms.
Thus the descended contraction becomes the specified complex
contraction.  Stein contraction and normality are preserved by these
algebraically closed extensions.

For a birational step, recover its positive model over $T$ as
\begin{equation}\label{glob:adjoint-algebra}
 S^+=\Proj_T\mathcal R,
 \qquad
 \mathcal R=\bigoplus_{m\geq0}f_*\cO_S(mbD),
\end{equation}
with $b>0$ sufficiently divisible.  Over $\C$, the effective exceptional
pullback difference \eqref{glob:step-difference} identifies these
pushforward sheaves with the corresponding sheaves on the positive
model, whose adjoint is relatively ample.  Hence the relative Proj is
exactly that model.  This includes a divisorial step for which the
positive model is $T$: the algebra is then the nonnegative section
algebra of the Cartier multiple $bD_T$, whose relative Proj is $T$.

For clarity, both the sheaf algebra and its finite generation descend.
On a normal variety over the algebraically closed field $F$, the sheaf
of a Weil divisor is the rank-one reflexive extension of its invertible
restriction on a regular open containing all codimension-one points.
Field extension is flat with geometrically regular fibres in
characteristic zero.  The codimension-two complement retains its
codimension, and the pullback sheaf retains reflexivity and the $S_2$
property.  It is therefore the same reflexive extension as the sheaf of
the extended divisor.  Multiplication is also compatible, since it
agrees on the regular open and its target is torsion-free.  The
canonical, boundary and nef-trace divisors defining $D$ are all defined
over $F$; the last assertion is checked by computing the trace on a
common resolution.

Choose $b$ clearing their rational coefficients and making the complex
extension of $bD$ Cartier.  The divisorial sheaf $\cO_S(bD)$ becomes
invertible after the faithfully flat field extension and so is
invertible already over $F$.  The same $b$ thus works downstairs.
Proper flat base change identifies every graded pushforward in
\eqref{glob:adjoint-algebra} and its multiplication after extension.
On a finite affine cover of $T_\C$, finite generation of the complex
algebra gives a common bound $d_0$ for the degrees of generators.
Downstairs, form the subalgebra generated by its coherent pieces of
degrees at most $d_0$.  In each degree its inclusion in $\mathcal R$
becomes surjective after faithful flat extension, so its cokernel is
zero.  Thus these finitely many coherent pieces generate $\mathcal R$,
proving finite generation without presupposing it.  Formation of
relative Proj commutes with extension, so it supplies the next normal
projective model and its birational map over $F$.  This inductively
descends the entire finite diagram, including its final fibre
contraction when present.

Finally extend the diagram from $F$ to $k$.  Projectivity, normality,
birationality, the absence of extraction on the positive side,
connected contractions and relative ampleness persist.  For a proper
contraction, connectedness here is also expressed by the equality of
the pushforward structure sheaf with the target structure sheaf, which
is preserved by flat base change.  Cartierness of the specified
multiples has already descended from $\C$ and is preserved over $k$.

We record the numerical facts needed for the endpoint.  Every line
bundle on an algebraically closed extension of a fixed projective
$F$-variety is numerically equivalent to one from $F$: spread the bundle
on a finite-type parameter scheme, specialize to an $F$-point, and use
degree constancy on its geometrically connected base change, as above.
Conversely, numerical triviality and nonnegativity of an $F$-divisor on
relative curves persist under extension.  Spread a vertical projective
curve, together with the point of the target containing its image, in
a flat projective family over a finite-type $F$-scheme.  A specialization
at an $F$-point is an effective vertical curve cycle with the same
intersection degrees.  The corresponding nonnegativity or vanishing
over $F$ therefore gives the same property for the original curve.
The reverse implication is tested on base changes of $F$-curves.
It follows that nefness and the relative numerical divisor spaces,
hence relative Picard ranks, are unchanged.  In the nef case the final
adjoint remains nef; in the Mori case its negative remains relatively
ample and the relative Picard number remains one.

Take common smooth resolutions of the descended finite steps, also
dominating the original model carrying the nef divisor, and extend them
to $k$.  Their effective pullback differences remain effective, with
the same positive coefficients along original contracted primes.
The calculation \eqref{glob:generalized-monotonicity} then gives weak
increase of every generalized log discrepancy over $k$, with strict
increase for each prime on the original $X$ contracted by the composite.
In particular the output is generalized lc.  Its boundary is the
pushforward of the original boundary and its nef part is the trace of
the original nef b-divisor. The same effective pullback differences show
that a nef endpoint makes the original adjoint pseudo-effective.
Conversely, pseudo-effectivity
of the original adjoint persists through the birational steps and rules
out an anti-ample restriction to a general positive-dimensional fibre.
Thus the sign of the original adjoint selects the stated alternative
over $k$ as well, completing Theorem~\ref{thm:main}.
\end{proof}

It remains to prove Theorem~\ref{thm:phase} and the local assertions on
which its proof depends.
